% Source: https://artofproblemsolving.com/community/c3h1216440_1st_reun_mock_middle_school_math_contest
% Stable identity: aops-topic-1216440-q19
By Heron's formula with semiperimeter $s=\dfrac{7+12+15}{2}=17$,

\begin{asy}
usepackage("amsmath");
size(5cm);
defaultpen(fontsize(10pt));
pair A=(0,0),B=(15,0),C=(10.6,5.497);
draw(A--B--C--cycle,linewidth(0.9)); draw(C--(C.x,0),dashed);
label("$15$",(A+B)/2,S); label("$12$",(A+C)/2,NW); label("$7$",(B+C)/2,NE);
label("\mbox{}",(min(currentpicture,user=true)+max(currentpicture,user=true))/2,black); // keeps the text after the figure black when LaTeX typesets the labels inline
\end{asy}

\[\text{Area}=\sqrt{17(17-7)(17-12)(17-15)}=\sqrt{17\cdot10\cdot5\cdot2}=\sqrt{1700}=\boxed{10\sqrt{17}}.\]
